\chapter{Tropical spectra of complex spaces}
\index{tropical spectrum}\index{semivaluation!complex space}
\label{chap:tropical-spectra-complex}

Berkovich analytification begins with a normed ring.
A compact complex space has no global algebraic coordinate ring and no
non-Archimedean ground norm from which to start.
Mesquita-Piccione's construction replaces functions by the idempotent
semiring of coherent ideals and takes its tropical spectrum
\cite{MP24}.
The output is a compact space of semivaluations that is intrinsic to the
complex space.

This chapter develops only that foundation:
the semiring, its spectrum, topology, supports, centers, functoriality, and
the recovery of the usual trivially valued analytification in the proper
algebraic case.
The psh and finite-energy theories built on this space require additional
arguments and are reserved for later chapters.

The min-plus convention is chosen to match the preceding valuative
chapters.  A larger ideal has a smaller ideal value, multiplication adds
values, and \(\log|\mathfrak a|=-v(\mathfrak a)\).  Thus ideal sums become
maxima of logarithmic functions.  Switching to max-plus without changing
all three conventions would reverse both the topology-generating
coordinates and the psh order.

\section{The tropical target}

Put
\[
 \mathbf T_+:=([0,+\infty],\oplus,\odot),
 \qquad
 a\oplus b:=\min\{a,b\},
 \qquad
 a\odot b:=a+b.
\tag{58.1}\label{eq:ts-tropical-semiring}
\]
The conventions are
\[
                 a+(+\infty)=+\infty,
 \qquad
                 \min\{a,+\infty\}=a.
\]
Thus the additive zero is \(+\infty\), the multiplicative unit is \(0\),
and addition is idempotent.
Endow \([0,+\infty]\) with its compact order topology.

Let \(S\) be an idempotent commutative semiring.

\begin{definition}[Tropical spectrum]
\label{def:ts-tropical-spectrum}
The tropical spectrum of \(S\) is
\[
 \operatorname{tSpec}(S)
       :=\operatorname{Hom}_{\mathrm{Semiring}}(S,\mathbf T_+),
\tag{58.2}\label{eq:ts-tSpec}
\]
endowed with the topology of pointwise convergence.
\end{definition}

Equivalently, view it as a subspace of the product
\([0,+\infty]^S\).

\begin{remark}[Positive target]
\label{rem:ts-positive-target}
For a general semiring, \cite[Definition~1.2.1]{MP24} uses the larger
target \(((-\infty,+\infty],\min,+)\).  Characters of an ideal semiring are
automatically nonnegative, so the spectrum relevant to this chapter lands
in \(\mathbf T_+\).  We build that positivity into
\Cref{def:ts-tropical-spectrum}; no claim is made that the restricted
definition agrees with the general tropical spectrum for an arbitrary
idempotent semiring.
\end{remark}

\begin{proposition}\label{prop:ts-compact}
The space \(\operatorname{tSpec}(S)\) is compact Hausdorff.
\end{proposition}

\begin{proof}
The ambient product is compact Hausdorff by Tychonoff's theorem.
The equations
\[
\begin{aligned}
 \chi(0_S)&=+\infty,&
 \chi(1_S)&=0,\\
 \chi(a+b)&=\min\{\chi(a),\chi(b)\},&
 \chi(ab)&=\chi(a)+\chi(b)
\end{aligned}
\]
cut out a closed subset because the operations on
\([0,+\infty]\) are continuous in the compact order topology.
For extended addition, the only endpoint to check is \(+\infty\): if
\(a_i\to+\infty\) and \(b_i\geq0\), then
\(a_i+b_i\to+\infty\).  Hence multiplication in \(\mathbf T_+\) is
continuous even at the absorbing element.
\end{proof}

The order on an idempotent semiring is
\[
                         a\leq b\quad\Longleftrightarrow\quad a+b=b.
\]
Every tropical character reverses it:
\[
                a\leq b\quad\Longrightarrow\quad\chi(a)\geq\chi(b).
\tag{58.3}\label{eq:ts-order-reversal}
\]

\section{The ideal semiring of a ring}

Let \(R\) be a commutative ring and
\(\mathscr I(R)\) its set of finitely generated ideals.
Set
\[
             I\oplus J:=I+J,\qquad I\odot J:=IJ.
\tag{58.4}\label{eq:ts-ideal-operations}
\]
Its additive zero is \((0)\), its multiplicative unit is \(R\), and
\(I+I=I\).

A nonnegative semivaluation on \(R\) is a map
\[
                       v:R\longrightarrow[0,+\infty]
\]
such that
\[
\begin{aligned}
 v(0)&=+\infty,&v(1)&=0,\\
 v(fg)&=v(f)+v(g),&
 v(f+g)&\geq\min\{v(f),v(g)\}.
\end{aligned}
\tag{58.5}\label{eq:ts-semivaluation}
\]
It is automatically trivial on every unit.

\begin{theorem}\label{thm:ts-affine-homeomorphism}
The assignments
\[
\begin{aligned}
 v&\longmapsto
   \chi_v,\qquad
   \chi_v((f_1,\ldots,f_r))
       :=\min_i v(f_i),\\
 \chi&\longmapsto
   v_\chi,\qquad
   v_\chi(f):=\chi((f))
\end{aligned}
\tag{58.6}\label{eq:ts-affine-correspondence}
\]
are inverse homeomorphisms between the space of nonnegative
semivaluations on \(R\), with pointwise topology, and
\(\operatorname{tSpec}\mathscr I(R)\).
\end{theorem}

\begin{proof}
The value \(\chi_v(I)\) does not depend on the chosen generators.
Indeed, if \(g=\sum a_if_i\) and \(v(a_i)\geq0\), then
\[
                      v(g)\geq\min_i v(f_i).
\]
Applying this in both directions to two generating sets proves equality.
The identities
\[
\begin{aligned}
 \chi_v(I+J)&=\min\{\chi_v(I),\chi_v(J)\},\\
 \chi_v(IJ)&=\chi_v(I)+\chi_v(J)
\end{aligned}
\]
show that \(\chi_v\) is a tropical character.

Conversely, multiplicativity of principal ideals gives
\[
 v_\chi(fg)=\chi((f)(g))=v_\chi(f)+v_\chi(g).
\]
Since \((f+g)\subseteq(f,g)=(f)+(g)\), order reversal gives
\[
 v_\chi(f+g)\geq
 \chi((f,g))
 =\min\{v_\chi(f),v_\chi(g)\}.
\]
Thus \(v_\chi\) is a semivaluation.
Finally, if \(I=(f_1,\ldots,f_r)\), then
\[
 \chi(I)=\chi((f_1)+\cdots+(f_r))
        =\min_i\chi((f_i)).
\]
The maps are inverse.
Continuity in both directions follows because each ideal has finitely many
generators and the topology on each side is pointwise.
\end{proof}

\begin{example}[A one-parameter tropical family on the affine line]
Let \(R=\mathbf C[x]\).  For \(0\leq\alpha<+\infty\), define
\[
 v_\alpha\!\left(\sum_i a_ix^i\right)
       :=\min_{a_i\ne0}i\alpha.
\]
At the endpoint set \(v_\infty(f)=0\) when \(f(0)\ne0\), and
\(v_\infty(f)=+\infty\) when \(x\mid f\).  These are nonnegative
semivaluations.  Under \Cref{thm:ts-affine-homeomorphism}, their characters
satisfy, for positive integers \(r,s,m\),
\[
 \chi_\alpha((x^m))=m\alpha,
 \qquad
 \chi_\alpha((x^r,x^s))=\min\{r,s\}\alpha.
\]
Thus ideal multiplication becomes addition of slopes, while ideal sum takes
their minimum exactly as in \eqref{eq:ts-ideal-operations}.

For finite \(\alpha\), no nonzero polynomial has value \(+\infty\); at
\(\alpha=+\infty\), precisely the polynomials in \((x)\) do.  The extended
ideal function and its bounded truncation are
\[
       \log|(x^m)|(v_\alpha)=-m\alpha,
       \qquad
       \max\{\log|(x^m)|(v_\alpha),-\ell\}
          =\max\{-m\alpha,-\ell\}.
\]
The second expression is real-valued and continuous even at
\(\alpha=+\infty\), illustrating why the truncation in
\eqref{eq:ts-ideal-truncation} belongs to the PL lattice whereas the bare
ideal function may take \(-\infty\).
\end{example}

If \(R\) is a finitely generated algebra over a trivially valued field, the
nonnegative condition is equivalent to boundedness of the multiplicative
seminorm
\[
                              |f|_v:=\rme^{-v(f)}
\]
by the trivial norm.
Thus \Cref{thm:ts-affine-homeomorphism} identifies the affine tropical
spectrum with the bounded, or centered, part of the usual Berkovich
spectrum.

\section{The ideal semiring of a complex space}

Let \(X\) be a complex space.
Denote by \(\mathscr I_X\) the set of globally defined ideal sheaves that
are locally of finite type.
For complex spaces these ideals are coherent by Oka's coherence theorem.
Define
\[
             \mathfrak a\oplus\mathfrak b
                   :=\mathfrak a+\mathfrak b,
 \qquad
             \mathfrak a\odot\mathfrak b
                   :=\mathfrak a\mathfrak b.
\tag{58.7}\label{eq:ts-sheaf-semiring}
\]

\begin{definition}[Intrinsic semivaluation space]
\label{def:ts-complex-beth}
The intrinsic non-Archimedean space of \(X\) is
\[
                     X^\bethan
                     :=\operatorname{tSpec}(\mathscr I_X).
\tag{58.8}\label{eq:ts-beth}
\]
A point is written \(v\), with
\[
 v(\mathfrak a)\in[0,+\infty].
\]
\end{definition}

By \Cref{prop:ts-compact}, \(X^\bethan\) is compact Hausdorff.
Its topology is the coarsest topology for which every evaluation
\[
             v\longmapsto v(\mathfrak a)
\tag{58.9}\label{eq:ts-evaluation}
\]
is continuous.
No choice of coordinates, embedding, polarization, K\"ahler class, or
resolution enters the definition.

\begin{remark}[Reducedness and irreducibility]
\label{rem:ts-generality}
Definition \eqref{eq:ts-beth} makes sense for an arbitrary complex space.
Statements that identify a point with a valuation of a function field, or
that speak of a full-support stratum, will assume the relevant irreducible
reduced subspace.
We do not silently replace a nonreduced space by its reduction.
\end{remark}

\section{Tropical PL functions}

Every coherent ideal defines an extended-real function
\[
 \log|\mathfrak a|:X^\bethan\longrightarrow[-\infty,0],
 \qquad
 \log|\mathfrak a|(v):=-v(\mathfrak a).
\tag{58.10}\label{eq:ts-ideal-function}
\]
The semiring laws become
\[
\begin{aligned}
 \log|\mathfrak a\mathfrak b|
      &=\log|\mathfrak a|+\log|\mathfrak b|,\\
 \log|\mathfrak a+\mathfrak b|
      &=\max\{\log|\mathfrak a|,\log|\mathfrak b|\}.
\end{aligned}
\tag{58.11}\label{eq:ts-ideal-function-laws}
\]
These equalities retain \(-\infty\) in the usual extended-real sense.

\begin{definition}[Tropical PL lattice]
\label{def:ts-pl-lattice}
Let \(\PL^+(X^\bethan)\) be the set of all \emph{real-valued} functions
of the form
\[
 v\longmapsto
 \frac1m\max_{1\leq i\leq N}
       \{\log|\mathfrak a_i|(v)+k_i\},
 \qquad
 m\in\mathbf Z_{>0},\quad k_i\in\mathbf Z,
\tag{58.12}\label{eq:ts-pl-positive}
\]
where the \(\mathfrak a_i\) are coherent ideals.  The real-valued
requirement is part of the definition; it excludes a maximum that is
\(-\infty\) at some support stratum.  Define
\(\PL_{\mathbf Q}(X^\bethan)\) to be the
\(\mathbf Q\)-vector space generated by \(\PL^+(X^\bethan)\).
\end{definition}

The class \(\PL^+\) is stable under finite maxima after clearing
denominators.  It contains every bounded truncation
\[
          \max\{\log|\mathfrak a|,-\ell\},
          \qquad \ell\in\mathbf Q_{>0},
\tag{58.13}\label{eq:ts-ideal-truncation}
\]
because the second term is obtained from the unit ideal.  Thus the
definition does not require a bare ideal function to be finite on every
support stratum.

\begin{proposition}\label{prop:ts-pl-density}
The lattice \(\PL_{\mathbf Q}(X^\bethan)\) separates points of
\(X^\bethan\).  Its real scalar extension is uniformly dense in
\(C(X^\bethan)\).
\end{proposition}

\begin{proof}
Distinct characters \(v,w\) differ on some
\(\mathfrak a\in\mathscr I_X\).  After interchanging them, suppose
\(w(\mathfrak a)<v(\mathfrak a)\), and choose
\(\ell\in\mathbf Q\) strictly between these two values (with any finite
\(\ell>w(\mathfrak a)\) if \(v(\mathfrak a)=+\infty\)).
Then \eqref{eq:ts-ideal-truncation} has different values at \(v\) and \(w\).
Hence \(\PL_{\mathbf Q}\) separates points.
The lattice Stone--Weierstrass theorem on the compact Hausdorff space
\(X^\bethan\) gives density.
\end{proof}

The truncation in this proof is a genuine PL function, but it does not
redefine the extended ideal function
\eqref{eq:ts-ideal-function}, whose value \(-\infty\) records a genuine
support condition.

\section{Support and center}

Assume throughout this section that $X$ is a compact reduced
irreducible complex space.
We use the strong Noetherian property in
\cite[Lemma~2.1.1 and Remark~2.1.5]{MPTropicalV2}: an increasing
sequence of coherent ideal sheaves is locally stationary.
Compactness makes such a sequence stationary globally, by taking
a finite cover and the largest of the corresponding indices.
Consequently any sum of coherent ideal sheaves is a finite sum.
Indeed, otherwise one could successively add an ideal not contained
in the preceding finite sum and obtain a strictly increasing sequence.
This uses local stationarity of sheaves, not just Noetherianity of
an individual stalk.

We also specify the meaning of prime used in this section, following
\cite[Notations and Conventions, p.~6]{MPTropicalV2}.
A proper coherent ideal $\mathfrak p$ is \emph{globally prime} if,
for coherent ideals $\mathfrak a,\mathfrak b$ on $X$,
\[
 \mathfrak a\mathfrak b\subseteq\mathfrak p
 \quad\Longrightarrow\quad
 \mathfrak a\subseteq\mathfrak p
 \ \text{or}\ \mathfrak b\subseteq\mathfrak p.
\tag{58.13a}\label{eq:ts-global-prime}
\]
This is equivalent to $\mathfrak p$ being the radical ideal of a
nonempty irreducible analytic subspace.
To see the forward implication, the coherent radical
$\mathfrak q=\sqrt{\mathfrak p}$ satisfies
$\mathfrak q^N\subseteq\mathfrak p$ for some $N$: local finite
generation gives such a power near each point, and compactness
gives a common exponent. Repeated use of
\eqref{eq:ts-global-prime} gives $\mathfrak q\subseteq\mathfrak p$,
so $\mathfrak p$ is radical.
Its zero set has finitely many irreducible components by compactness.
If there were at least two, take the ideals of one component and
of the union of the remaining components.
Their product is contained in $\mathfrak p$, but neither ideal is,
contradicting \eqref{eq:ts-global-prime}.
Conversely, if $Y$ is irreducible and
$\mathfrak a\mathfrak b\subseteq\mathfrak i_Y$, then
$Y\subseteq V(\mathfrak a)\cup V(\mathfrak b)$.
Irreducibility and the analytic Nullstellensatz give one of the
two ideal inclusions in \eqref{eq:ts-global-prime}.
This global condition does not assert that every proper stalk
$\mathfrak p_x$ is a prime ideal of $\mathcal O_{X,x}$.

For $v\in X^\bethan$, define
\[
\begin{aligned}
 \mathfrak p_s(v)
   &:=\sum_{v(\mathfrak a)=+\infty}\mathfrak a,\\
 \mathfrak p_c(v)
   &:=\sum_{v(\mathfrak a)>0}\mathfrak a.
\end{aligned}
\tag{58.14}\label{eq:ts-support-center-ideals}
\]
Both families contain the zero ideal. By the finite-sum argument,
these are coherent ideals and their values are respectively
$+\infty$ and strictly positive.

\begin{proposition}\label{prop:ts-prime-support-center}
The ideals $\mathfrak p_s(v)$ and $\mathfrak p_c(v)$ are globally
prime, hence define nonempty reduced irreducible analytic subspaces.
Moreover,
\[
                         \mathfrak p_s(v)\subseteq\mathfrak p_c(v).
\tag{58.15}\label{eq:ts-prime-inclusion}
\]
\end{proposition}

\begin{proof}
Their positive values exclude the unit ideal, whose value is zero.
If $\mathfrak a\mathfrak b\subseteq\mathfrak p_s(v)$, order
reversal and multiplicativity give
\[
 +\infty=v(\mathfrak p_s(v))
       \leq v(\mathfrak a)+v(\mathfrak b).
\]
One summand is infinite, so its ideal is contained in
$\mathfrak p_s(v)$ by the definition of that sum.
For $\mathfrak p_c(v)$, the same reasoning gives
$v(\mathfrak a)+v(\mathfrak b)>0$; nonnegativity forces at
least one summand to be positive. Thus both ideals satisfy
\eqref{eq:ts-global-prime}. The inclusion follows because an
infinite value is positive.
\end{proof}

\begin{samepage}
\begin{definition}
\label{def:ts-support-center}
The \emph{support} and \emph{center} of $v$ are
\[
 \Supp(v):=V(\mathfrak p_s(v)),
 \qquad
 Z(v):=V(\mathfrak p_c(v)),
\tag{58.16}\label{eq:ts-support-center}
\]
with their reduced analytic structures.
\end{definition}
\end{samepage}

Both are irreducible analytic subspaces, and
\[
                              Z(v)\subseteq\Supp(v).
\tag{58.17}\label{eq:ts-center-in-support}
\]

\begin{example}
\label{ex:ts-node-global-prime}
Consider the projective nodal cubic
\[
             C=\{y^2z=x^2(x+z)\}\subset\mathbf P^2.
\]
The polynomial $y^2-x^2(x+1)$ is irreducible over $\mathbf C[x,y]$,
since $x^2(x+1)$ is not a square in $\mathbf C(x)$.
Thus $C$ is an irreducible projective curve, and its trivial
ideal character has $\mathfrak p_s=\mathfrak p_c=0$ and
$\Supp(v_{\triv})=Z(v_{\triv})=C$.
At $p=[0:0:1]$, however, the analytic local ring is
\[
 \mathcal O_{C,p}
   =\mathbf C\{x,y\}/\bigl(y^2-x^2(1+x)\bigr).
\]
The convergent square root of $1+x$ gives two nonzero germs
$y-x\sqrt{1+x}$ and $y+x\sqrt{1+x}$ whose product is zero.
Hence the zero ideal is globally prime on $C$, while its stalk
at the node is not prime. Global support strata retain this
irreducible curve even though its germ has two branches.
\end{example}

\begin{proposition}\label{prop:ts-support-stratification}
There is a disjoint decomposition
\[
                X^\bethan
                =\coprod_{Y\subseteq X}Y^{\val},
\tag{58.18}\label{eq:ts-support-strata}
\]
over nonempty reduced irreducible analytic subspaces $Y$.
Here $Y^{\val}$ is the space of ideal semivaluations on $Y$
that are finite on every nonzero coherent ideal, mapped to
$X^\bethan$ by $\mathfrak a\mapsto w(\mathfrak a\mathcal O_Y)$.
Its image consists exactly of the characters with support $Y$.
\end{proposition}

\begin{proof}
Let $\Supp(v)=Y$ and $\mathfrak i_Y=\mathfrak p_s(v)$.
For a coherent ideal $\mathfrak b$ on $Y$, take its inverse
image $\widetilde{\mathfrak b}$ under
$\mathcal O_X\to\mathcal O_Y$ and set
$w(\mathfrak b)=v(\widetilde{\mathfrak b})$.
The inverse image is coherent and contains $\mathfrak i_Y$.
Any other lift has the same value, because adjoining
$\mathfrak i_Y$ does not change a value:
\[
 v(\mathfrak a+\mathfrak i_Y)
       =\min\{v(\mathfrak a),+\infty\}=v(\mathfrak a).
\]
Sums of lifts lift sums, and the product of two lifts becomes
the inverse image of the product after adjoining $\mathfrak i_Y$.
Thus $w$ satisfies both semiring laws, with $w(0)=+\infty$
and $w(\mathcal O_Y)=0$.
If $w(\mathfrak b)=+\infty$, its inverse image is contained
in $\mathfrak p_s(v)=\mathfrak i_Y$, so $\mathfrak b=0$.
It follows that $w\in Y^{\val}$, and the displayed map in the
statement recovers $v$.

Conversely, for $w\in Y^{\val}$, the character
$v(\mathfrak a)=w(\mathfrak a\mathcal O_Y)$ has infinite value
exactly when $\mathfrak a\subseteq\mathfrak i_Y$.
Consequently $\mathfrak p_s(v)=\mathfrak i_Y$ and its support
is $Y$. These constructions are inverse, and uniqueness of the
support makes the decomposition disjoint.
\end{proof}

Let $p\in X$ and let $\mathfrak m_p$ denote its coherent ideal sheaf.

\begin{proposition}\label{prop:ts-point-criteria}
For $v\in X^\bethan$,
\[
\begin{aligned}
 Z(v)=\{p\}
       &\quad\Longleftrightarrow\quad v(\mathfrak m_p)>0,\\
 \Supp(v)=\{p\}
       &\quad\Longleftrightarrow\quad v(\mathfrak m_p)=+\infty.
\end{aligned}
\tag{58.19}\label{eq:ts-point-criteria}
\]
\end{proposition}

\begin{proof}
The ideal $\mathfrak m_p$ is maximal among proper coherent
ideals on $X$. If $v(\mathfrak m_p)>0$, its inclusion in the
proper ideal $\mathfrak p_c(v)$ therefore gives equality, and
$Z(v)=\{p\}$. Conversely, a radical coherent ideal with zero
set $\{p\}$ is exactly $\mathfrak m_p$, so this center condition
gives $v(\mathfrak m_p)=v(\mathfrak p_c(v))>0$.
The same argument for $\mathfrak p_s(v)$ proves the support criterion.
\end{proof}

\begin{proposition}\label{prop:ts-local-global-point}
The assignment
\[
 w\longmapsto v_w,\qquad
 v_w(\mathfrak a):=w(\mathfrak a_p),
\]
is a bijection from nonnegative semivaluations of
$\mathcal O_{X,p}$ that are positive on its maximal ideal to
the points of $X^\bethan$ centered at $p$.
\end{proposition}

\begin{proof}
Finite stalk generators show that $v_w$ satisfies the semiring
laws and $v_w(\mathfrak m_p)>0$.
Thus \Cref{prop:ts-point-criteria} gives its center.

Conversely, suppose $Z(v)=\{p\}$ and put
$c=v(\mathfrak m_p)\in(0,+\infty]$.
Write $A=\mathcal O_{X,p}$ and $\mathfrak m=\mathfrak m_{X,p}$.
For $f\in A$ and $k\geq1$, the ideal
$(f)+\mathfrak m^k$ determines a coherent ideal
$\mathfrak a_k(f)$ on $X$ that is the unit ideal off $p$.
One construction takes the inverse image of the ideal generated
by the class of $f$ modulo $\mathfrak m^k$ in
$\mathcal O_X/\mathfrak m_p^k$.
This quotient is supported at $p$, so the construction depends
only on the germ and requires no global extension of $f$.
The ideals decrease in $k$, and hence their values increase.
Define
\[
       w(f):=\lim_{k\to\infty}v(\mathfrak a_k(f)).
\tag{58.19a}\label{eq:ts-local-recovery}
\]

The unit has value zero, whereas
$v(\mathfrak a_k(0))=kc\to+\infty$, so $w(0)=+\infty$.
Here $k(+\infty)=+\infty$, which includes $c=+\infty$.
The exact ideal inclusions
\[
 \mathfrak a_{2k}(fg)
   \subseteq\mathfrak a_k(f)\mathfrak a_k(g)
   \subseteq\mathfrak a_k(fg)
\]
hold at $p$ by expanding
$((f)+\mathfrak m^k)((g)+\mathfrak m^k)$, and off $p$
all three ideals are the unit ideal.
Order reversal and multiplicativity give
\[
 v(\mathfrak a_{2k}(fg))
  \geq v(\mathfrak a_k(f))+v(\mathfrak a_k(g))
  \geq v(\mathfrak a_k(fg)).
\]
Taking increasing limits proves $w(fg)=w(f)+w(g)$.
There is no subtraction of infinite values.
Similarly,
$\mathfrak a_k(f+g)\subseteq\mathfrak a_k(f)+\mathfrak a_k(g)$
gives the semivaluation inequality for sums.
If $f\in\mathfrak m$, then
$\mathfrak a_k(f)\subseteq\mathfrak m_p$, so $w(f)\geq c>0$.

The reconstruction also has an exact finite-level form.
For $\ell\geq k$,
$\mathfrak a_\ell(f)+\mathfrak m_p^k=\mathfrak a_k(f)$.
Applying $v$ and letting $\ell\to\infty$ yields
\[
              v(\mathfrak a_k(f))=\min\{w(f),kc\}.
\tag{58.19b}\label{eq:ts-local-finite-truncation}
\]
For a global coherent ideal $\mathfrak a$, choose finite
generators $f_1,\ldots,f_r$ of $\mathfrak a_p$, taking $f_1=0$
if that stalk is zero. Then $\mathfrak a+\mathfrak m_p^k=\sum_j\mathfrak a_k(f_j)$.
Finite minima commute with these increasing limits, giving
\[
 \begin{split}
 w(\mathfrak a_p)
 &=\lim_{k\to\infty}v(\mathfrak a+\mathfrak m_p^k)\\
 &=\lim_{k\to\infty}\min\{v(\mathfrak a),kc\}
   =v(\mathfrak a).
 \end{split}
\]
Thus $v_w=v$.
Starting with $w$ instead, the same truncated-ideal formula gives
$v_w(\mathfrak a_k(f))=\min\{w(f),kw(\mathfrak m)\}$,
whose limit is $w(f)$ because $w(\mathfrak m)>0$.
This proves the reverse composite and the bijection.
It completes the reconstruction in
\cite[Example~2.1.4, p.~12]{MPTropicalV2}.
\end{proof}

\section{Functoriality}

Let \(f:X\to Y\) be a holomorphic map.
Pullback of ideal sheaves means the image in \(\cO_X\) of
\[
                    f^{-1}\mathfrak a\otimes_{f^{-1}\cO_Y}\cO_X.
\]
It respects ideal sums, products, zero, and the unit.
Therefore it is a semiring homomorphism
\[
                         f^*:\mathscr I_Y\longrightarrow\mathscr I_X.
\]

\begin{definition}
\label{def:ts-functorial-map}
Define
\[
 f^\bethan:X^\bethan\longrightarrow Y^\bethan,
 \qquad
 f^\bethan(v)(\mathfrak a):=v(f^*\mathfrak a).
\tag{58.20}\label{eq:ts-functorial-map}
\]
\end{definition}

\begin{proposition}\label{prop:ts-functoriality}
The map \(f^\bethan\) is continuous,
\[
 (g\circ f)^\bethan=g^\bethan\circ f^\bethan,
 \qquad
 (\mathrm{id}_X)^\bethan=\mathrm{id}_{X^\bethan}.
\tag{58.21}\label{eq:ts-functoriality}
\]
If \(f\) is a closed embedding, then \(f^\bethan\) is a topological
embedding.
\end{proposition}

\begin{proof}
For every \(\mathfrak a\in\mathscr I_Y\), evaluation of
\(f^\bethan(v)\) at \(\mathfrak a\) is evaluation of \(v\) at
\(f^*\mathfrak a\), hence is continuous.
The composition identities follow from functoriality of ideal pullback.

For a closed embedding, every coherent ideal on \(X\) is the quotient of a
coherent ideal on \(Y\).
Thus \(f^*:\mathscr I_Y\to\mathscr I_X\) is surjective, so precomposition on
tropical characters is injective.
It is a continuous injection from a compact space to a Hausdorff space,
hence a topological embedding.
\end{proof}

An open embedding does not in general induce an open embedding on tropical
spectra.
In particular, \(U^\bethan\) is compact even when \(U\subset X\) is a
proper noncompact open set.
The construction is global in coherent ideals, not a naive topological
gluing of open Berkovich charts.

\begin{proposition}\label{prop:ts-bimeromorphic-valuations}
Let \(f:X'\to X\) be a proper bimeromorphic morphism of reduced irreducible
compact complex spaces.
Then \(f^\bethan\) induces a bijection between the strata of valuations with
full support \(X'\) and \(X\).
\end{proposition}

\begin{proof}
First, a nonzero ideal on \(X\) has nonzero pullback to \(X'\): on any open
set where the ideal is nonzero, bimeromorphicity prevents its inverse image
from being contained in the proper zero locus.  Thus \(f^\bethan\) sends a
full-support character to a full-support character.

We import the relative-ampleness construction of
\cite[Proposition~2.1.6]{MP24}, but record its formulas.  Suppose first
that \(X'\) admits an anti-effective \(f\)-ample divisor \(E\).  For every
coherent ideal \(\mathfrak j\) on \(X'\), relative global generation gives
an integer \(m_{\mathfrak j}\) and ideals
\(\mathfrak i_{\mathfrak j},\mathfrak k\) on \(X\) such that
\[
\begin{aligned}
 f^{-1}\mathfrak i_{\mathfrak j}\cdot\cO_{X'}
   &=\mathfrak j\cdot\cO_{X'}(E)^{m_{\mathfrak j}},\\
 f^{-1}\mathfrak k\cdot\cO_{X'}
   &=\cO_{X'}(E).
\end{aligned}
\]
If \(v\in X^{\val}\), define
\[
 \widetilde v(\mathfrak j)
 :=v(\mathfrak i_{\mathfrak j})
      -m_{\mathfrak j}v(\mathfrak k).
\]
All quantities are finite, which is precisely where full support is used.
Passing two ideals to a common sufficiently large twist and using the
semiring laws shows that this value is independent of the choices and that
\(\widetilde v\) is a tropical character.  The displayed pullback
identities then give
\[
 f^\bethan(\widetilde v)=v.
\]
They also recover the value of every \(\mathfrak j\) from its pushforward
data, proving injectivity in this case.

For a general proper bimeromorphic \(f\), the resolution and negativity
argument in \cite[Proposition~2.1.6]{MP24} produces a smooth space \(Z\)
dominating both \(X'\) and \(X\), with the morphism \(Z\to X\) a sequence
of blowups whose anti-effective exceptional divisor is relatively ample.
The special case gives bijections on full-support strata for the resulting
morphisms.  In the commutative triangle, surjectivity for \(Z\to X'\) and
bijectivity for \(Z\to X\) force both injectivity and surjectivity for
\(f^\bethan:X'^{\val}\to X^{\val}\).
\end{proof}

We assert only a bijection on the full-support strata.
A statement about all support strata, or a homeomorphism without additional
hypotheses, would be stronger and is not used here.

\section{Recovery of projective algebraic analytification}

Let \(X\) now be a complex algebraic variety and
\(X^{\mathrm{an,triv}}\) its Berkovich analytification with respect to the
trivial norm on \(\mathbf C\).

\begin{theorem}\label{thm:ts-algebraic-recovery}
If \(X\) is proper, there is a canonical functorial homeomorphism
\[
                         X^\bethan
                    \xrightarrow{\ \sim\ }
                         X^{\mathrm{an,triv}}.
\tag{58.22}\label{eq:ts-algebraic-recovery}
\]
\end{theorem}

\begin{proof}
GAGA identifies algebraic coherent ideals with analytic coherent ideals when
\(X\) is proper, so passing to the associated complex space does not
enlarge or shrink the ideal semiring.  For the algebraic locally ringed
space, \cite[Example~1.3.2(1)]{MP24} identifies the tropical spectrum with
the subspace of the Berkovich analytification consisting of semivaluations
centered on \(X\).  Properness makes every semivaluation centered, by the
valuative criterion of properness.  The construction is functorial, and
the pointwise topologies agree, proving the claimed homeomorphism.

For a nonproper algebraic scheme the cited result still applies to the
algebraic ideal semiring.  We do not identify that semiring with the
semiring of analytic coherent ideals without a GAGA hypothesis.
\end{proof}

Under \eqref{eq:ts-algebraic-recovery}, the ideal function
\(\log|\mathfrak a|\) is the usual model ideal function
\[
                              v\longmapsto-v(\mathfrak a),
\]
and the functorial map \eqref{eq:ts-functorial-map} is ordinary Berkovich
analytification.
Thus the tropical construction genuinely extends, rather than replaces, the
projective algebraic space used in the preceding chapter.

\begin{proposition}
\label{prop:ts-multiplier-bridge}
Let \(X\) be smooth and projective, let \(L\) be a pseudo-effective
\(\mathbf Q\)-line bundle, and let
\(\varphi\in\PSH(X,c_1(L,h))\).  Under
\eqref{eq:ts-algebraic-recovery},
\[
 \varphi^{\An}(v)
 =\inf_{j\geq0}2^{-j}
       \log|\cI(2^j\varphi)|(v).
\tag{58.23}\label{eq:ts-multiplier-bridge}
\]
For a DXZ metric \(\Gamma\), choose any K\"ahler perturbation \(\omega\).
Its divisorial analytification is therefore
\[
 \Gamma^{\An}(v)
 =\sup_{\tau<\Gamma_{\max}}
   \left\{\tau+\inf_{j\geq0}2^{-j}
     \log|\cI(2^j\Gamma^{\theta+\omega}_\tau)|(v)\right\}
 \quad(v\in X^{\divis}).
\tag{58.24}\label{eq:ts-test-curve-bridge}
\]
\end{proposition}

\begin{proof}
By \eqref{eq:ts-ideal-function},
\[
 2^{-j}\log|\cI(2^j\varphi)|(v)
 =-2^{-j}v(\cI(2^j\varphi)).
\]
Taking the infimum and applying the dyadic multiplier-ideal formula
\eqref{eq:va-dyadic} proves
\eqref{eq:ts-multiplier-bridge}.  Substitute that equality level by level
in the Legendre formula \eqref{eq:comp-legendre} to obtain
\eqref{eq:ts-test-curve-bridge}.  The second identity is kept on divisorial
valuations because that is the proved raw Legendre domain before the BJ
extension theorem is invoked.  It is independent of \(\omega\), since a
DXZ transition preserves every integer multiplier ideal.  In particular,
no unperturbed component \(\Gamma^\theta\) has been introduced in the
zero-volume case.
\end{proof}

The order of the infimum and supremum in
\eqref{eq:ts-test-curve-bridge} is part of the construction.  Fekete's
lemma first completes the multiplier-ideal data of one complex level; only
then does Legendre duality range over \(\tau\).  No minimax interchange is
used.

\section{Boundary of this chapter}

The results above are the tropical-spectrum foundations of \cite{MP24}.
They provide a compact semivaluation space for complex spaces and recover the
usual projective trivially valued analytification.
They do not yet define a psh cone, Monge--Amp\`ere energy, entropy, or a
comparison with the test-curve theory.

In particular, \Cref{thm:ts-algebraic-recovery} identifies the underlying
spaces and ideal evaluations.
It does not by itself identify the Mesquita-Piccione psh class with the
Darvas--Xia--Zhang or Boucksom--Jonsson psh class.
The missing positivity comparison is supplied later by
\Cref{thm:pseff-projective-full-mp-bj}: for a connected smooth projective
\(X\) and a real ample N\'eron--Severi class \(\alpha\), the full MP and BJ
psh cones coincide under \eqref{eq:ts-algebraic-recovery}, with no
finite-energy hypothesis.  For an ample \(\mathbf Q\)-class, composing
this equality with the DXZ--BJ comparison gives the full DXZ--MP order
isomorphism on all semivaluations.  This does not cover a nonprojective
compact K\"ahler manifold, a transcendental class outside the real
N\'eron--Severi space, or a non-nef class, and it does not identify
Monge--Amp\`ere operators merely from equality of their psh domains.

\begin{remark}[Primary-source boundary]
The ideal-semiring tropical spectrum, its PL lattice, support
stratification, the local--global bijection at a point, functoriality,
bimeromorphic full-support statement, and proper algebraic recovery are from
\cite[Sections~1.2--2.1]{MP24}.  The affine ideal-character calculation is
the elementary semiring argument of
\Cref{thm:ts-affine-homeomorphism}; GAGA is the additional input in
\Cref{thm:ts-algebraic-recovery}.  \Cref{prop:ts-multiplier-bridge} only
rewrites the already proved DXZ
multiplier-ideal formula in these coordinates.  Neither this proposition nor
the underlying-space homeomorphism alone asserts equality of psh cones; the
additional model-positivity and decreasing-completion argument is
\Cref{thm:pseff-projective-full-mp-bj}.
\end{remark}
